\subsection{训练样本与目标}

三种 LoRA 适配器（\texttt{LoRA\_code}、\texttt{LoRA\_cot\_zs}、\texttt{LoRA\_cot\_fs(k)}）共用 \texttt{KodCode\_train.jsonl}，但在 assistant 回复的构造与 loss 计算口径上有所区分，以匹配各自对应的推理场景。

\paragraph{（1）三类监督目标}
同一训练集经不同脚本处理后形成三种监督目标（见表~\ref{tab:ch4-lora-target}），支撑 H4/H7/H9 三条训练路径；具体脚本与指针文件对应 4.1 节表~\ref{tab:ch4-matrix}。
\begin{table}[!htb]
  \centering
  \caption{三种 LoRA 适配器的监督目标与训练 prompt 配置}
  \label{tab:ch4-lora-target}
  \footnotesize
  \begin{tabular}{llll}
    \toprule
    适配器 & assistant 回复内容 & 训练 prompt few-shot & 产出脚本 \\
    \midrule
    \texttt{LoRA\_code}       & \texttt{solution}（纯代码块） & 无 & \texttt{train/train.py} \\
    \texttt{LoRA\_cot\_zs}    & \texttt{thought\_step} + \texttt{solution} & 0 例 & \texttt{train\_cot.py --num-few-shots 0} \\
    \texttt{LoRA\_cot\_fs($k$)} & \texttt{thought\_step} + \texttt{solution} & $k$ 例 & \texttt{train\_cot.py --num-few-shots k} \\
    \bottomrule
  \end{tabular}
\end{table}

\paragraph{（2）Chat 样本结构}
每条训练样本一律封装为 OpenAI chat messages 三元组：system 固化角色，user 承载「题目 + 函数签名 + 指令」，assistant 给出监督目标。与基线 Direct 的唯一差别在于 CoT 组的 user instruction 末尾追加 \texttt{Think step by step, then output the Python code.}，以及 assistant 采用固定的「推理段 + $\backslash$n$\backslash$n + \verb|```python ... ```|」格式——两者共同保证训练 prompt 与 4.4 节推理侧第二阶段的上下文分布严格一致。

\paragraph{（3）Loss 掩码策略}
LoRA 微调只应对 assistant 段求 loss，否则会在 system/user 上消耗容量、干扰指令跟随。本文采用的掩码构造思路如下：
\begin{itemize}
    \item \textbf{完整序列}：\texttt{apply\_chat\_template(messages, tokenize=True)} 得到完整 token 序列 $T$；
    \item \textbf{非 assistant 前缀}：\texttt{apply\_chat\_template(messages[:-1], add\_generation\_prompt=True)} 得到截至 assistant 起始标记的前缀 $P$；
    \item \textbf{掩码对齐}：比较 $T$ 与 $P$ 的公共前缀长度 $L$，令 \texttt{labels[:L] = -100}，其余位置保留原 token id；
    \item \textbf{长度约束}：\texttt{MAX\_LENGTH=1024} 为硬上限，超长样本直接截断（训练阶段只保留能完整容纳 assistant 的样本），避免因截断落在 assistant 中段导致的断句偏差。
\end{itemize}
该策略只对最后一条 assistant 段反向传播，既保持 chat template 对特殊 token 的原生解释，又避免手工拼字符串带来的 tokenizer 不一致。

\paragraph{（4）Few-shot 训练样本的构造}
\texttt{LoRA\_cot\_fs(k)} 独有环节：对每条目标样本 $q$，用 3.3.2 节定义的强标签集合派生 \texttt{strong\_tags}，在训练池中按
\[
\text{score}(q, q') = 2\cdot |\mathrm{strong\_tags}(q) \cap \mathrm{strong\_tags}(q')| + |\mathrm{tags}(q) \cap \mathrm{tags}(q')|
\]
评分候选，取得分最高的前 $\max(k, 4k)$ 条作为候选池，再用 \texttt{random.sample} 从中随机抽取 $k$ 条拼在 user 段之前，构成 \texttt{(题目, 推理, 代码)} 形式的示例序列。该构造与 4.4.2 节评测侧的检索函数共享同一评分与同一强标签集合，仅在随机性来源上存在差异（训练侧 \texttt{random.sample}，评测侧用 \texttt{seed = 2026 + problem\_idx * 1000 + sample\_idx} 做可复现抽样），以此保证 train/test 上下文统计一致而评测仍可稳定复现。
